\section*{Chapitre \ref{ch:mx:the-empirics-of-price-impact} --- Les études empiriques de l'impact de prix}

\begin{solution}{exo:mx:the-empirics-of-price-impact:1}
$0.8\times0.02\times\sqrt{0.01}=0.0016$, 16~points de base~; $0.8\times0.02\times\sqrt{0.04}=0.0032$, 32~points de base.
\end{solution}

\begin{solution}{exo:mx:the-empirics-of-price-impact:2}
Environ $0.6\times2.1\approx1.3$ \omterm{def:mx:the-limit-order-book:tick}{pas de cotation}. Les petits déséquilibres rencontrent les niveaux peu épais près du prix et le déplacent proportionnellement~; les grands atteignent des niveaux plus profonds et
réveillent les fondamentalistes, de sorte que la relation s'infléchit.
\end{solution}

\begin{solution}{exo:mx:the-empirics-of-price-impact:3}
Tout, à son erreur près~: $-1.5$ \omterm{def:mx:the-limit-order-book:tick}{pas de cotation}, avec une erreur type de 1.5. Les fondamentalistes échangent vers la valeur cachée, que le métaordre n'a pas
déplacée, et ramènent le prix en arrière.
\end{solution}

\begin{solution}{exo:mx:the-empirics-of-price-impact:4}
La dérive prévue vaut $a\sigma\sqrt T$ et $T$ est proportionnel à $Q/V$ à taux de participation fixe, si bien que la dérive est proportionnelle à $\sigma\sqrt{Q/V}$,
de même forme que l'impact~: elle s'ajoute au \omterm{def:mx:the-empirics-of-price-impact:prefactor}{préfacteur d'impact} ($0.8+a$ fois la part d'ordres temporisés, à peu près) et laisse la pente inchangée.
\end{solution}

\begin{solution}{exo:mx:the-empirics-of-price-impact:5}
Les ordres portant sur un même titre partagent sa volatilité, sa liquidité et ses séances de cotation, si bien que leurs erreurs sont corrélées~; rééchantillonner des ordres isolés les traite comme indépendants
et donne des intervalles trop étroits.
\end{solution}

\begin{solution}{exo:mx:the-empirics-of-price-impact:6}
Ses ordres ont été enregistrés sur un marché sans les nôtres~: ses fournisseurs de liquidité continuent d'afficher là où ils affichaient et ses preneurs de liquidité continuent de prendre ce qu'ils prenaient, quoi que
nous fassions. Un backtest sur données rejouées ne tient compte que du coût mécanique du carnet d'ordres consommé, sans réaction à celui-ci, et sous-estime donc le coût des grands ordres.
\end{solution}

\begin{solution}{exo:mx:the-empirics-of-price-impact:7}
Sans ordres temporisés, l'exposant tombe à 0.28 (préfacteur à exposant fixé 0.81)~; avec eux, 0.20 (1.24). Sur un dixième de journée, le bruit
$\sigma\sqrt T$ est le même pour tous les ordres et submerge les petits, dont les moyennes par classe sont dominées par le bruit~; la prévision des ordres temporisés, désormais identique
quelle que soit la taille, aplatit encore la courbe.
\end{solution}

\begin{solution}{exo:mx:the-empirics-of-price-impact:8}
L'exposant peut être juste alors que le niveau ne l'est pas~: dans la simulation, 30~\% des ordres temporisés par une prévision laissent l'exposant à 0.50 et font monter le préfacteur
de 62~\%. Vérifier le préfacteur sur des ordres sans prévision, ou retirer la prévision ordre par ordre.
\end{solution}

\begin{solution}{pb:mx:the-empirics-of-price-impact:1}
\textbf{1.} Fournisseurs de liquidité plaçant des ordres par rapport au carnet courant, ordres de bruit tirés à l'avance, fondamentalistes échangeant vers une valeur cachée~; les fournisseurs et les
fondamentalistes réagissent.
\textbf{2.} 0.27, 0.58, 1.12 et 2.34 \omterm{def:mx:the-limit-order-book:tick}{pas de cotation} pour 1, 2--3, 4--7 et 8+ lots~; après cinquante ordres 0.81, 0.23, 0.68 et 0.46.
\textbf{3.} Les fondamentalistes informés envoient des ordres d'un lot, et le prix se déplace vers la valeur à leur suite.
\textbf{4.} 2.1 \omterm{def:mx:the-limit-order-book:tick}{pas de cotation}.
\textbf{5.} La séance avec le métaordre moins la même séance sans lui, avec les mêmes ordres de bruit et les mêmes sauts de valeur~: le bruit exogène s'annule.
\textbf{6.} Aucune différence à aucun horizon~: les fournisseurs de liquidité rejoués n'apprennent jamais l'existence de l'ordre.
\textbf{7.} 6.3, 7.8 et $-1.5$ \omterm{def:mx:the-limit-order-book:tick}{pas de cotation}.
\textbf{8.} Les écarts de prix après le premier ordre enfant croissent avec les choix différents des fournisseurs~; des erreurs types de 1 à 3 \omterm{def:mx:the-limit-order-book:tick}{pas de cotation} face à des impacts de quelques
\omterm{def:mx:the-limit-order-book:tick}{pas de cotation}.
\textbf{9.} $I=Y\sigma(Q/V)^\delta$, $\delta\approx\tfrac12$~; 16 et 32 points de base.
\textbf{10.} Impact moyen sur la volatilité dans des classes de taille, un ajustement log-log sur les moyennes par classe, un bootstrap sur les titres~; les impacts isolés peuvent être négatifs.
\textbf{11.} $\delta=0.48$ (0.39 à 0.56), $Y=0.85$ (0.55 à 1.66)~; 0.81 à exposant fixé.
\textbf{12.} Un carnet moyen en V s'annulant au prix, et une racine carrée s'ajustant sur deux ordres de grandeur de taille, et un logarithme sur cinq.
\textbf{13.} Le mouvement de prix attendu par le donneur d'ordre, compté comme impact.
\textbf{14.} $\delta=0.50$ (0.44 à 0.55), $Y=1.42$ (1.01 à 2.04), 1.31 à exposant fixé.
\textbf{15.} L'horizon de la prévision est la durée, qui croît avec la taille, si bien que la dérive a la forme de l'impact.
\textbf{16.} Retrancher la prévision de chaque ordre temporisé~; la régression de l'impact sur la prévision donne 1.54, car les ordres plus grands portent les deux effets.
\textbf{17.} L'exposant tombe à 0.28 sans ordres temporisés et à 0.20 avec eux~: le bruit et la prévision indépendante de la taille aplatissent la courbe.
\textbf{18.} Sélection des ordres exécutés, métaordres qui se chevauchent, contextes rejoués (et bruit).
\textbf{19.} \emph{Résultat nommé}~: sur 2\,000 métaordres simulés de loi vraie 0.8 et un demi, l'ajustement donne $\delta=0.48$ (0.39 à 0.56) et $Y=0.85$ (0.55
à 1.66)~; avec 30~\% des ordres temporisés par une prévision, $\delta=0.50$ et $Y=1.42$, le préfacteur trop élevé de 62~\%, l'exposant restant inchangé.
\textbf{20.} Une racine carrée avec le préfacteur ajusté sur des ordres sans prévision, ou avec la prévision de chaque ordre retirée.
\end{solution}

\begin{solution}{iq:mx:the-empirics-of-price-impact:1}
$I=Y\sigma\sqrt{Q/V}$~: le mouvement de prix d'un métaordre en fraction du prix, $Y$ de l'ordre de un, $\sigma$ la volatilité quotidienne, $Q$ la taille de l'ordre et $V$
volume quotidien~; elle vaut en moyenne, pour des ordres qui ne sont qu'une fraction du volume quotidien, quel que soit le calendrier d'exécution dans de larges limites.
\iqlookfor{Chaque terme~; ce qui est moyenné~; le domaine de validité.}
\end{solution}

\begin{solution}{iq:mx:the-empirics-of-price-impact:2}
Du début à la dernière exécution par rapport au milieu à l'arrivée, normalisé par la volatilité, réparti en classes selon la taille rapportée au volume, ajusté en logarithmes, regroupé par titre et par jour. Biais~:
prévisions de calendrier, sélection des ordres exécutés, chevauchement avec d'autres ordres, bruit et horizon choisi.
\iqlookfor{Normalisation~; répartition en classes~; regroupement~; \omterm{def:mx:the-empirics-of-price-impact:alpha}{contamination par l'alpha}.}
\end{solution}

\begin{solution}{iq:mx:the-empirics-of-price-impact:3}
La liquidité près du prix est faible et croît à mesure qu'on s'en éloigne (une liquidité latente en V)~: consommer $Q$ exige un mouvement dont le carré croît avec $Q$. Par ailleurs,
les grands ordres sont découpés et le carnet se reconstitue entre les ordres enfants.
\iqlookfor{Liquidité latente~; découpage~; reconstitution.}
\end{solution}

\begin{solution}{iq:mx:the-empirics-of-price-impact:4}
Le marché rejoué ne peut pas réagir~: seul le carnet d'ordres consommé par l'ordre est facturé, et personne ne fait bouger les prix contre lui. Il faut utiliser un modèle d'impact ou un
simulateur réactif.
\iqlookfor{Limite du rejeu~; une correction.}
\end{solution}

\begin{solution}{iq:mx:the-empirics-of-price-impact:5}
Une partie du pic revient après l'exécution et une partie demeure~; l'ampleur dépend du caractère informé du flux, de la vitesse de retour de la liquidité et de ce qui s'échange par ailleurs.
Dans la simulation, tout est revenu, car rien n'y apprend de l'ordre~; sur données réelles, une part permanente est courante.
\iqlookfor{Retour~; information~; horizon.}
\end{solution}

\begin{solution}{iq:mx:the-empirics-of-price-impact:6}
Pas nécessairement~: leurs ordres sont temporisés par des prévisions, et la dérive prévue est comptée comme impact. Comparer après retrait de la prévision, ou sur des ordres envoyés
sans prévision~; les ordres temporisés de la simulation affichaient un préfacteur supérieur de 62~\% à exécution identique.
\iqlookfor{\omterm{def:mx:the-empirics-of-price-impact:alpha}{Contamination par l'alpha}~; une comparaison équitable.}
\end{solution}
